\documentclass[11pt]{article}
\usepackage[a4paper,margin=25mm]{geometry}
\usepackage{amsmath,amssymb,amsthm}
\usepackage[hidelinks]{hyperref}
\emergencystretch=2em
\newtheorem{theorem}{Theorem}
\newcommand{\R}{\mathbb R}
\title{Equal Equilibria and Different Relaxation Spectra:\\
A Proof of Conjecture 00000006640}
\author{GitHub: gaochengzhecpu}
\date{}
\begin{document}
\maketitle
\begin{abstract}
Two symmetric two-state master equations, with jump rates one and two,
have the same static solution space and the same unique stationary
probability vector. Their generator spectra are $\{0,-2\}$ and
$\{0,-4\}$, and their propagation spectra differ at every positive time.
The construction, stochastic semigroups, differential equations, exact
spectra, and existential claim are verified in Lean using actual real
matrices and derivatives.
\end{abstract}

\paragraph{Reading of the source.}
Both language versions ask for an explicit pair of master equations
with identical static solution sets but different time spectra.
For a finite-state continuous-time master equation, write probability
vectors as columns and use
\[
 p'(t)=Qp(t),\qquad p_i(t)\geq0,\qquad \sum_i p_i(t)=1.
\]
The generator has nonnegative off-diagonal entries and zero column sums.
Static probability solutions are $\{p:Qp=0,\ p_i\geq0,\ \sum_i p_i=1\}$;
without probability normalization the static solutions form $\ker Q$.
We establish equality in both senses. We establish different spectra
both for the generator and for the propagator at every fixed $t>0$.
The generators and propagators are real symmetric and obey detailed
balance for the same positive probability vector, so the example also
has the usual time-reversal symmetry. Time is measured in the same
units in both systems; the source imposes no quotient by time rescaling.

\begin{theorem}\label{thm:main}
For $r>0$ define the two-state generator and probability vector
\[
 Q_r=\begin{pmatrix}-r&r\\r&-r\end{pmatrix},\qquad
 \pi=\begin{pmatrix}1/2\\1/2\end{pmatrix}.
\]
The systems $p'=Q_1p$ and $p'=Q_2p$ satisfy:
\begin{enumerate}
\item They are valid reversible master equations with the common
unique stationary probability vector $\pi$, and
$\ker Q_1=\ker Q_2=\{(a,a)^{\mathsf T}:a\in\R\}$.
\item Their generator spectra are $\{0,-2\}$ and $\{0,-4\}$.
\item Their stochastic fundamental solutions are
\begin{equation}\label{eq:P}
 P_r(t)=\frac12\begin{pmatrix}
 1+e^{-2rt}&1-e^{-2rt}\\
 1-e^{-2rt}&1+e^{-2rt}
 \end{pmatrix},
\end{equation}
with $P_r(0)=I$, $P_r(s+t)=P_r(s)P_r(t)$, and
$P_r'(t)=Q_rP_r(t)$. For every initial probability vector $p_0$,
$p(t)=P_r(t)p_0$ remains a probability vector for $t\geq0$ and solves
the master equation with $p(0)=p_0$.
\item For each $t>0$, the propagator spectra
$\{1,e^{-2t}\}$ and $\{1,e^{-4t}\}$ differ.
\end{enumerate}
Thus the source's requested separation exists, including for symmetric
reversible systems.
\end{theorem}

\begin{proof}
For $r>0$ the off-diagonal generator entries equal $r$ and both column
sums vanish. The matrix is symmetric and
$(Q_r)_{ij}\pi_j=(Q_r)_{ji}\pi_i$, proving detailed balance.
For a column $p=(x,y)^{\mathsf T}$,
\[
 Q_rp=r(y-x,x-y)^{\mathsf T}.
\]
It vanishes exactly when $x=y$. Imposing $x+y=1$ yields
$x=y=1/2$, a nonnegative vector. This proves both static assertions.

The characteristic determinant is
\[
 \det(\mu I-Q_r)=(\mu+r)^2-r^2=\mu(\mu+2r).
\]
For a finite real matrix the real algebraic spectrum consists exactly
of values for which this determinant vanishes, so
$\sigma(Q_r)=\{0,-2r\}$. The two chosen spectra are distinct, since
$-2$ belongs to the first but not the second. Equivalently, the constant
mode $(1,1)^{\mathsf T}$ has eigenvalue zero and the imbalance mode
$(1,-1)^{\mathsf T}$ has eigenvalue $-2r$.

In \eqref{eq:P}, $0<e^{-2rt}\leq1$ for $r,t\geq0$, so the entries
are nonnegative and the columns sum to one. Direct substitution gives
$P_r(0)=I$. The identity
$e^{-2r(s+t)}=e^{-2rs}e^{-2rt}$ gives the semigroup law by matrix
multiplication. Differentiating the entries gives
\[
 P_r'(t)=r e^{-2rt}
 \begin{pmatrix}-1&1\\1&-1\end{pmatrix}=Q_rP_r(t).
\]
These identities hold for all real $s,t$; stochasticity is used only
for nonnegative time. Multiplication by a fixed initial column gives
$p'=Q_rp$ and the stated initial condition. Nonnegativity and column
sums preserve the probability constraints. Also $P_r(t)\pi=\pi$.

Put $a=e^{-2rt}$. A second determinant calculation gives
\[
 \det(\mu I-P_r(t))=(\mu-1)(\mu-a),
 \qquad \sigma(P_r(t))=\{1,e^{-2rt}\}.
\]
For $t>0$, strict monotonicity of the real exponential gives
$0<e^{-4t}<e^{-2t}<1$. Therefore $e^{-2t}$ belongs to the first
propagator spectrum but not the second. The two decay rates are
different even though their static states coincide.
\end{proof}

\section*{Formal verification and scope}
The file \texttt{lean/Main.lean} uses Lean 4.19.0 and Mathlib commit
\begin{center}\texttt{c44e0c8ee63ca166450922a373c7409c5d26b00b}.\end{center}
The state space is \texttt{Fin 2}; generators and propagators are actual
$2\times2$ matrices over $\R$. The probability predicate includes both
nonnegativity and normalization. The static predicate is defined using
actual matrix-vector multiplication.

The formalization verifies the generator conditions, Hermitian
symmetry, detailed balance, the unique normalized static vector, and
equality of the full unnormalized kernels. The exact spectra use
Mathlib's algebraic \texttt{spectrum}, not a list defined to be the
claimed eigenvalues. They are obtained by reducing noninvertibility
to the determinant and proving the determinant factorizations.
The propagators satisfy the initial condition and semigroup law.
They preserve nonnegativity and total probability. Their derivatives
are actual \texttt{HasDerivAt} statements. For every initial vector,
the trajectory satisfies the master equation in each coordinate.
This finite-dimensional existence claim does not require a probability
path construction.

The final theorem is
\begin{center}\texttt{Conjecture6640.conjecture\_6640\_true}.\end{center}
It proves \texttt{SeparationClaim}: there exist two generators, two
stochastic fundamental solutions, and a common detailed-balance
probability vector; the normalized static sets and full kernels
coincide, the generator spectra differ, and the propagator spectra
differ at every positive time. The actual mathematical objects and
the complete existential statement are thus present in Lean.
The proof does not invoke a general matrix-exponential or ODE
uniqueness theorem; it verifies the displayed explicit fundamental
solutions directly. All mathematical assertions required by the
theorem are formalized. The displayed eigenvector explanation is
an additional paper explanation of the formally proved determinant
calculation.

The audited main and bridge theorems use only \texttt{propext},
\texttt{Classical.choice}, and \texttt{Quot.sound}. This submission
was developed by a delegated AI agent with a self-review; the
coordinating agent reviews it separately. No external independent
review is claimed.

\section*{Source and reproduction}
The exact English and Chinese source appears in \texttt{SOURCE.md},
from upstream commit
\begin{center}\texttt{b355e42906e6eaf52bc7e9d005e0ef02d8aa7400}.\end{center}
Its hash and retrieval information are in
\begin{center}\texttt{verification/SOURCE\_PROVENANCE.md}.\end{center}
The argument is elementary and self-contained. The standard
$Q$-matrix and master-equation framework is discussed, for example,
in J.~R.~Norris, \emph{Markov Chains}, Chapter 2,
\href{https://www.statslab.cam.ac.uk/~jrn10/Markov/}{author's book page}.
The column-vector convention used here is explicitly specified above.

With the pinned dependencies available, run in the \texttt{lean}
directory:
\begin{center}\texttt{lake env lean -DwarningAsError=true Main.lean}\end{center}
and \texttt{lake build} for a fresh build. Compile \texttt{main.tex}
with Tectonic or a compatible LaTeX engine.
\end{document}
